\input{preamble.tex}

\title{\huge{First assignment}}
\author{\LARGE{Thobias Høivik}}
\date{}

\begin{document}
\maketitle

This is the first assignment for MAT242. As the 
course materials, announcements, etc. are in English, I 
will be writing this in English. That is my preference anyways. 
Hopefully the use of the word "write-up" doesn't necessitate 
this being super-duper structured and dull, but that I am free 
to share my exploratory thought process. Still I will strive 
to be as rigorous as possible, of course.  

Note: I am submitting one day late as I thought the deadline was 
the $7$th. My mistake, I have been away at a conference 
and things got a bit hectic. 


\section*{The problem}
If I understood the problem correctly, I am to show 
that the sort of "canonical" way of defining a topology 
on the cartesian product of some sets $X_i$ is the only way 
of doing it when we want to satisfy that universal property, 
whatever that means!

I have heard enough waffling on Category Theory to know that 
this term is relevant there and I was immediately confused by 
the "there is only one way to do it" line. A cursory 
google search on "product topology vs. box topology" immediately 
brought up how this so-called "universal property" fails to 
be satisfied by the box topology when looking at the cartesian 
product of infinitely many $X_i$. 

\section*{The universal property}

This property is stated for the cartesian product 
of two sets so I will begin there. 
I will denote the projection $\operatorname{pr}_i$ 
by $\pi_i : X_1 \times X_2 \to X_i$ because it looks cooler and 
writing \verb|\operatorname| a bunch is tedious.

Let $X_1,X_2$ be sets. The cartesian product has the property 
that for any function $f : Z \to X_1 \times X_2$, 
$f$ is uniquely defined by the two projections 
\[
    f_1 := \pi_1 f : Z \to X_1 \ \text{and} \ 
    f_2 := \pi_2 f : Z \to X_2, 
\]
where $\pi_i : X_1 \times X_2 \to X_i$ is the projection 
$(x_1,x_2) \overset{\pi_i}{\mapsto} x_i$ for $i = 1,2$.

This essentially says something quite trivial which is that 
"to specify where $z$ goes in $X_1 \times X_2$, specify its 
$X_1$-coordinate and its $X_2$-coordinate". 

We would want the same to hold for topological spaces. 
If $X_1,X_2, Z$ are topological spaces, then a function 
$f : Z \to X_1 \times X_2$ is continuous if and only if 
$f_1$ and $f_2$ are. 

There is probably some nice way to formulate the above 
as requiring the some diagram commutes or something, but I 
don't know Category Theory as of right now. 

What we now want to do is to show that this requirement forces 
a unique topology on $X_1\times X_2$ called the product topology. 

\section*{The product topology}
So we are essentially given that this is the product topology. 
That means that what we should try to do now is to prove that 
for arbitrary spaces $X_1,X_2,Z$, the product topology has the 
desired property. 

\begin{defn*}[Product topology]
    Let $X_1$ and $X_2$ (for notational consistency) be topological  
    spaces, meaning there exist some $\tau_1 \subseteq \mathcal P(X_1)$
    and $\tau_2 \subseteq \mathcal P(X_2)$ such that 
    $(X_1,\tau_1)$ and $(X_2, \tau_2)$ are topological spaces. 

    Consider 
    \[
        \mathcal B 
        = \{U \times V \bigm| U \subseteq X_1 \ \text{open}, \ 
        \text{and} \ V \subseteq X_2 \ \text{open}\}.
    \]
    The product topology on $X_1 \times X_2$ is the 
    topology having $\mathcal B$ as a basis. 

    So an arbitrary open set of $X_1\times X_2$ is a union 
    of $U\times V$: 
    \[
        W = \bigcup_{\alpha \in A} U_\alpha \times V_\alpha. 
    \]
\end{defn*}
Munkres also gives an equivalent definition using projections, 
very exciting. 

Now suppose are given some continuous function 
$f : Z \to X_1 \times X_2$, meaning that 
for every open set $W \in X_1 \times X_2$ we have that 
\[
f^{-1}(W) = \{z \in Z \bigm| f(z) \in W\}
\]
is open in $Z$. 

Now we should think back to what this means under the universal 
property of the cartesian product, since this must necessarily 
hold for $f$.  
Let $f_i := \pi_if : Z \times X_i, \ (i = 1,2)$. Then 
$f(z) = (f_1(z), f_2(z))$. The idea now becomes quite clear. 
We assume $f$ is continuous, we know the projections are continuous 
(though we should probably prove so) and continuity is preserved 
under composition (maybe we ought to prove that aswell or maybe 
it won't be necessary) so 
$f$ continuous $\Rightarrow$ $f_1,f_2$ continuous. Then we must prove 
the converse which isn't that bad either. 

\begin{lem*}
    Let $X_1,X_2,Z$ be topological spaces and 
    endow $X_1 \times X_2$ with the product topology. If 
    \[
        f : Z \to X_1 \times X_2
    \]
    is continuous, then 
    \[
        f_1 = \pi_1 \circ f, \ f_2 = \pi_2 \circ f
    \]
    are continuous. 
\end{lem*}
\begin{proof}
    Suppose $f : Z \to X_1 \times X_2$ is continous. 
    Then, for any open set 
    in $X_1 \times X_2$, 
    \[
        f^{-1}(W) = \{z \in Z \bigm| f(z) \in W\}
    \]
    is open in $Z$. 

    Take $f_1$. Suppose $U \subseteq X_1$ is open. Then 
    \begin{align*}
        f_1^{-1}(U) &= (\pi_1 \circ f)^{-1}(U) \\ 
                    &= f^{-1}(\pi_1^{-1}(U)). 
    \end{align*}
    Then since 
    \[
        \pi_1^{-1}(U) = U \times X_2,
    \]
    and $U$ is open in $X_1$ by assumption, and $X_2$ is 
    open in itself, we have that $U \times X_2$ is open 
    in the product topology. 

    Therefore since $f$ is continuous, we have that 
    \[
        f_1^{-1}(U) = f^{-1}(U\times X_2) 
    \]
    is open in $Z$. 

    The proof of continuity for $f_2$ is effectively identical. 
\end{proof}
Perhaps it would have been good enough to remark that 
$f_i$ is the composition of continuous functions, but who knows. 
\begin{lem*}
    Let $X_1,X_2,Z$ be topological spaces and again 
    give $X_1\times X_2$ the product topology. If 
    \[
        f_1 = \pi_1 \circ f, \ f_2 = \pi_2 \circ f
    \]
    are continuous then 
    \[
        f(z) = (f_1(z), f_2(z))
    \]
    is continuous. 
\end{lem*}
\begin{proof}
    Suppose $W \subseteq X_1 \times X_2$ is open. 

    We want to show that 
    \[
        f^{-1}(W) 
    \]
    is open. 

    Recall that 
    \[
        f^{-1}(W) = f^{-1}\left(\bigcup_{\alpha \in A}U_\alpha 
        \times V_\alpha\right),
    \]
    where $U_\alpha, V_\alpha$ are open in $X_1, X_2$, respectively, 
    and for all $\alpha$.

    Recall also that 
    \begin{align*}
        f^{-1}\left(\bigcup_{\alpha} U_\alpha \times V_\alpha\right)
        &= \bigcup_{\alpha} f^{-1}(U_\alpha \times V_\alpha) \\ 
        &= \bigcup_{\alpha} f_1^{-1}(U_\alpha) \cap f_2^{-1}(V_\alpha),
    \end{align*}
    since 
    \begin{align*}
        z\in f^{-1}(U_\alpha\times V_\alpha)
            &\iff f(z)\in U_\alpha\times V_\alpha\\
            &\iff (f_1(z),f_2(z))\in U_\alpha\times V_\alpha\\
            &\iff f_1(z)\in U_\alpha\text{ and }f_2(z)\in V_\alpha\\
            &\iff z\in f_1^{-1}(U_\alpha)\cap f_2^{-1}(V_\alpha).
    \end{align*}
    Thus the preimage of an open set in $X_1 \times X_2$ is 
    the union of preimages (under $f_1,f_2$) of 
    open sets in $X_1$ and $X_2$. Since $f_1,f_2$ are continuous 
    by assumption, 
    $f^{-1}(W)$ is a union of open sets in $Z$ and therefore open.
\end{proof}

\begin{thm*}
    For topological spaces $X_1,X_2,Z$ where we once again 
    endow $X_1\times X_2$ with the product topology, 
    \[
        f : Z \to X_1 \times X_2
    \]
    is continuous if and only if 
    \[
        \pi_1 \circ f, \ \pi_2 \circ f
    \]
    are continuous. 
\end{thm*}
\begin{proof}
    Apply the previous two lemma. 
\end{proof}

If Category Theory has a more compact way of showing 
that the product topology indeed has this property I shall 
work tirelessly against those who label it as "abstract nonesense". 

\section*{It has to be the product topology!}
Again we are still working with simple 
cartesian products of only two topological spaces. 

Now what should the intuition be here, I wonder? 
Say we have some spaces $X_1, X_2$ and some topology on 
$X_1 \times X_2$. Why must it be that for any $Z$ and 
any 
functions $f : Z \to X_1 \times X_2$, 
$f$ cont. iff $f_1,f_2$ cont. must necessarily imply 
that $X_1\times X_2$ has the product topology? 

Well, this condition must hold for every $Z$ and $f$. So let 
$Z = X_1\times X_2$ with the product topology (which we 
now knows satisfies the nice property itself), and let 
$f = \operatorname{id}$. 

From here we can get a simple argument! 

\begin{thm*}
    Let $\tau$ be a topology on $X_1 \times X_2$. Suppose that 
    for every topological space $Z$ and every 
    $f : Z \to X_1 \times X_2$, $f$ is continuous iff 
    $\pi_1 f, \pi_2 f$ are continuous. Then 
    $\tau = \tau_{\mathrm{prod}}$. 
\end{thm*}
\begin{proof}
    Let $\tau_{\mathrm{prod}}$ denote the product topology on
    $X_1\times X_2$.

    We first show that
    \[
        \tau \subseteq \tau_{\mathrm{prod}}.
    \]
    Consider
    \[
        Z=(X_1\times X_2,\tau_{\mathrm{prod}})
    \]
    and the identity map
    \[
        \operatorname{id}:
        (X_1\times X_2,\tau_{\mathrm{prod}})
        \to
        (X_1\times X_2,\tau).
    \]
    Its coordinate maps are
    \[
        \pi_1\circ\operatorname{id}=\pi_1,
        \qquad
        \pi_2\circ\operatorname{id}=\pi_2.
    \]
    The projections $\pi_1$ and $\pi_2$ are continuous with respect
    to the product topology. Hence, by the assumed property of $\tau$,
    the identity map is continuous.

    Therefore every $\tau$-open set is
    $\tau_{\mathrm{prod}}$-open, and so
    \[
        \tau\subseteq\tau_{\mathrm{prod}}.
    \]

    We now prove the reverse inclusion. Consider
    \[
        Z=(X_1\times X_2,\tau)
    \]
    and the identity map
    \[
        \operatorname{id}:
        (X_1\times X_2,\tau)
        \to
        (X_1\times X_2,\tau_{\mathrm{prod}}).
    \]
    By applying the assumed property of $\tau$ to the identity map
    \[
        \operatorname{id}:
        (X_1\times X_2,\tau)
        \to
        (X_1\times X_2,\tau),
    \]
    we see that the projections
    \[
        \pi_1:(X_1\times X_2,\tau)\to X_1,
        \qquad
        \pi_2:(X_1\times X_2,\tau)\to X_2
    \]
    are continuous.

    Since the product topology satisfies the universal property,
    continuity of these two coordinate maps implies that
    \[
        \operatorname{id}:
        (X_1\times X_2,\tau)
        \to
        (X_1\times X_2,\tau_{\mathrm{prod}})
    \]
    is continuous. Hence
    \[
        \tau_{\mathrm{prod}}\subseteq\tau.
    \]

    Therefore
    \[
        \tau=\tau_{\mathrm{prod}}.
    \]
\end{proof}
Very cool! 

After this I was thinking back to the box topology thing which, 
to be honest, I haven't looked into that much. I was curious though 
to what I found about the box topology not satisfying this property 
when the cartesian product is infinite. After a quick search I found 
out that this topology is apparently the same as the product topology 
for finite cartesian product so that turned out to be relatively 
uninteresting. 

\end{document}

